% !TEX root = ../main.tex

\chapter{はじめに}
\label{chp:first}

\section{研究背景}
\label{sec:paragraph}

人は他者と影響を与え合いながら生きており，その関わりの中で価値観や行動，自己認識は少しずつ変化していく．
日常生活における会話や，何気ない振る舞いを目にすることなど，さまざまな他者との関わりが，
自分自身の考え方や行動に影響を与えている．しかし，この変化は多くの場合，無意識のうちに生じるものであり，
当人がその渦中で自覚することは難しい．\\
また，他者から受ける影響は必ずしも直接的なものとは限らない．誰かから言葉をかけられたり，
何かを教えられたりすることで考え方が変化する場合もあれば，直接影響を受けていなくとも，
他者同士が関わり合う様子を見続けるうちに，気づかないまま自分自身の価値観や行動が変化している場合もある．
このような変化は，身体的な変化とは異なり，目で見て確かめることができない．そのため，
「なぜ自分はこのような考え方をするようになったのか」，「これまでの誰との関わりが今の自分を作り上げたのか」を
辿ることは困難である．\\
さらに，他者から影響を受けて自分が変化するということは，自分自身もまた，
自覚のないまま誰かに影響を与えている可能性を意味する．何かを与えたつもりがなくても，自分の言葉や行動，
あるいは何気ない振る舞いが，誰かの考え方や行動に影響を与えているかもしれない．
しかし，自分が他者に与えた影響は，自分が他者から受けた影響以上にとらえることが難しい．\\
このように，自己は他者との関わりの中で形成され，新たな関わりによってさらに変化していくものだと考えられる．
そして，私たちは他者から影響を受けて変化するだけでなく，自らもまた他者の変化に関わっている可能性がある．
しかし，このような自己と他者の間に生じる影響や自己変容は，目に見える形で現れるものではなく，
日常の中でその過程を実感することは難しい．\\
以上のことから，本研究では他者との関わりの中で生じる自己の変化と，
その過程における自己と他者の相互関係について着目する．


\section{研究目的}
\label{sec:enum}

本研究の目的は，他者との関わりに応じて身体映像が互いに変化するインタラクティブ・インスタレーションを制作し，
自己と他者が影響し合う中で生じる自己変容の過程を体験的に問い直すことである．\\
ここでいう自己変容とは，経験や他者とのやりとりによって形成された自己が，
新たな関わりを通じてさらに変化していく過程を指す．なお，この変化は，自己と他者のどちらか一方にだけ生じる
変化ではなく，互いの関わりの中で生じる相互的な変化として捉える．\\
そこで，参加者の身体を映像として提示し，他者との距離や接触に応じて，それぞれの身体映像が変化する
インタラクションを設計する．自己の内面的な変化そのものは直接見ることはできないため，
本研究ではその変化を身体イメージの映像的な歪みとして表現し，参加者が自身の身体を通して知覚できる形で提示する．
